\documentclass[11pt]{article}

% ---------------------------------------------------------------- packages
\usepackage[margin=1in]{geometry}
\usepackage{amsmath, amssymb, mathtools}
\usepackage{bm}                 % bold symbols for vectors/matrices
\usepackage{siunitx}            % \qty{9.8}{\meter\per\second\squared}
\usepackage{booktabs}           % \toprule \midrule \bottomrule
\usepackage{graphicx}
\usepackage{microtype}
\usepackage[hidelinks]{hyperref}
\usepackage{cleveref}

\sisetup{per-mode = symbol}

% ---------------------------------------------------------------- macros
% Vectors and matrices
\newcommand{\vb}[1]{\bm{#1}}                       % vector:  \vb{v}
\newcommand{\mat}[1]{\bm{#1}}                      % matrix:  \mat{R}
\newcommand{\T}{^{\mathsf{T}}}                     % transpose
\newcommand{\skewm}[1]{\left[#1\right]_{\times}}   % skew-symmetric matrix

% Frames -- superscript/subscript labels
\newcommand{\fI}{\mathcal{I}}                      % inertial (NED)
\newcommand{\fB}{\mathcal{B}}                      % body (FRD)
\newcommand{\inI}[1]{{}^{\fI}\!#1}                 % quantity in inertial frame
\newcommand{\inB}[1]{{}^{\fB}\!#1}                 % quantity in body frame

% Common quantities
\newcommand{\Rbi}{\mat{R}}                         % body -> inertial rotation
\newcommand{\quat}{\vb{q}}
\newcommand{\qmul}{\otimes}                        % quaternion product
\newcommand{\eIII}{\vb{e}_3}                       % (0,0,1)

\newcommand{\ra}[1]{\renewcommand{\arraystretch}{#1}}
\newcommand{\degree}{\ensuremath{^\circ\,}}
\newcommand{\overbar}[1]{\mkern 1.5mu\overline{\mkern-1.5mu#1\mkern-1.5mu}\mkern 1.5mu}
\newcommand{\f}[2]{\frac{#1}{#2}}
\newcommand{\mb}[1]{\mathbf{#1}}
\newcommand{\tr}[1]{\mathrm{Tr}\left({#1}\right)}
\newcommand{\mbg}[1]{\boldsymbol{\mathbf{#1}}}
\DeclarePairedDelimiter{\ceil}{\lceil}{\rceil}
\DeclareMathAlphabet\mathbfcal{OMS}{cmsy}{b}{n}
\renewcommand{\d}{\mathop{}\!\mathrm{d}} % total derivative
\newcommand{\p}{\partial}
\renewcommand{\Re}{\mathrm{Re}}  % Real part (upright)
\renewcommand{\Im}{\mathrm{Im}}  % Imaginary part (upright)

\title{Quadrotor GNC Project}
\author{Max Howell}
\date{\today}

\begin{document}
\maketitle

\section{Introduction}

Intro placeholder


\section{Model and Derivation of EOM}

In this project, we consider the gz\_x500 quadcopter model used in Gazebo as shown in Figure 1. % TODO: add quadcopter figure and FBD next to it.
As can be seen, this quadcopter is in an X-configuration with motors 0-3 with the NED inertial frame $\mathcal{N}$ and the FRD body frame $\mathcal{B}$.
Motors 0 and 1 rotate counter clock-wise (CCW) and motors 2 and 3 rotate clock-wise (CW).
The geometry of the quadcopter is defined by radial vectors $\mb{r}_i$ which give the vector from the center of mass to each motor $i$ and the inertia is defined by the inertia matrix $\mb{J} = \mathrm{diag}(I_{xx}, I_{yy}, I_{zz})$.

The FBD shows the forces acting on the quadcopter: the thrust produced by each motor $\mb{T}_i = [0, 0, -T_i]$, the gravitational force $m\mb{g}$, where $m$ is the mass of the quadrotor and $\mb{g} = [0, 0, g]$ with $g$ the gravitational acceleration, and the rotor drag force $\mb{H}_i = -k_d |\omega_i| \mb{v}_{i,\perp}$, with $\omega_i$ being the rotational rate of the rotor and $\mb{v}_{i,\perp}$ being the perpendicular velocity.
% TODO: define $k_d$ (rotor drag coefficient) and give the expression for
% $\mb{v}_{i,\perp}$ -- the component of the rotor velocity perpendicular to the
% rotor axis, i.e. projected into the rotor plane.
There is also a reaction moment $\mb{M}_i$ produced by each spinning rotor as well as the moments produced by the forces.
% TODO: give the expression for $\mb{M}_i$ (reaction moment about $z_b$, with
% sign set by the rotor spin direction).

We derive the equations of motion using the Newton--Euler method and quaternions to define the rotation of the quadcopter.
The inertial position of the quadcopter is given by $\mb{p} = [p_x, p_y, p_z]^\top$ and the inertial velocity ${\mb{v}} = \dot{\mb{p}}$.
The quaternion state of the quadcopter defining the rotation is given by $\mb{q} = [q_0, q_1, q_2, q_3]^\top$, with $q_0$ collecting $\mathrm{cos}(\theta/2)$ and $q_{1:3}$ collecting $\mb{n}\mathrm{sin}(\theta/2)$, with $\theta$ being the angle of rotation and $\mb{n}$ the axis of rotation.
The angular rate of the quadcopter is given by $\boldsymbol{\omega} = [\omega_x, \omega_y, \omega_z]^\top$, with $\omega_x, \omega_y, \omega_z$ being the roll, pitch, and yaw rates in $\mathcal{B}$, respectively.
Collecting the full state and defining it as $\mb{x}$ gives
\begin{equation}
\mb{x} \coloneqq \begin{bmatrix}
\mb{p} \\ \mb{v} \\ \mb{q} \\ \boldsymbol{\omega}
\end{bmatrix}
=
\begin{bmatrix}
p_x \\ p_y \\ p_z \\ v_x \\ v_y \\ v_z \\ q_0 \\ q_1 \\ q_2 \\ q_3 \\ \omega_x \\ \omega_y \\ \omega_z
\end{bmatrix}
,
\end{equation}
which defines the 13-state model.

Summing the translational dynamics by $\sum{\mb{F}} = m\dot{\mb{v}}$ gives
\begin{equation}
\mb{R}(\sum{\mb{T_i}} + \sum{\mb{H_i}}) + m\mb{g} = m\dot{\mb{v}},
\end{equation}
where $\mb{R} \coloneqq \mb{R}(\mb{q})$ is the $3\times 3$ rotational matrix satisfying $\mb{v} = \mb{R}\mb{v}_b$, mapping from the body frame to the inertial frame.

Similarly, the rotational dynamics via $\sum{\boldsymbol{\tau}} = \mb{J}\dot{\boldsymbol{\omega}} + (\boldsymbol{\omega} \times \mb{J}\boldsymbol{\omega})$ gives
\begin{equation}
\sum(\mb{r}_i \times \mb{T}_i) + \sum{(\mb{r}_i \times \mb{H}_i)} + \sum{\mb{M}_i} = \mb{J}\dot{\boldsymbol{\omega}} + (\boldsymbol{\omega} \times \mb{J}\boldsymbol{\omega}),
\end{equation}
which can be re-arranged to solve for $\dot{\boldsymbol{\omega}}$.
The attitude kinematics comes from the quaternion derivative equation $\dot{\mb{q}} = 1/2 \mb{q} \otimes \begin{bmatrix}0 & \boldsymbol{\omega} \end{bmatrix}$.

Defining the control input as the thrust magnitude produced by each motor $\mb{u} = [T_0, T_1, T_2, T_3]^\top$, we wish to derive an allocation matrix $\mb{B}$ mapping from $\mb{u}$ to the total thrust as well as the roll, pitch, and yaw torques on the body.
The total thrust is simply a summation of the individual thrust from each motor $T_{\mathrm{tot}} = \sum{T_i}$.
The roll $\tau_x$ and pitch $\tau_y$ torques acting on the quadcopter body come from the cross product $(\mb{r}_i \times \mb{T}_i) = \begin{bmatrix} -r_{i,y} T_i & r_{i,x} T_i & 0 \end{bmatrix}$, where the torque produced by the thrust in the yaw direction is zero since it is perpendicular to the axis of thrust.
The yaw torque $\tau_z$ comes from the reaction moment produced by each motor $\mb{M}_i = \begin{bmatrix} 0 & 0 & \sigma_i k_M T_i \end{bmatrix}$, where $\sigma_{0,1} = 1$ for the CCW spinning rotors and $\sigma_{2,3} = -1$ for the CW spinning rotors, and $k_M$ is the moment constant for the rotor.
Collecting all the terms into a single mapping gives the allocation matrix
\begin{equation}
\mb{B} = 
\begin{bmatrix}
1 & 1 & 1 & 1\\
-r_{0,y}&  -r_{1,y}&  -r_{2,y} & -r_{3,y}\\
r_{0,x}&  r_{1,x}&  r_{2,x} & r_{3,x}\\
\sigma_0 k_M & \sigma_1 k_M & \sigma_2 k_M & \sigma_3 k_M
\end{bmatrix},
\end{equation}
which gives the mapping $\begin{bmatrix}T_{\mathrm{tot}} & \tau_x & \tau_y & \tau_z \end{bmatrix}^\top = \mb{B}\mb{u}$.

Re-writing the translational and rotational equations as a function of $\mb{B}$ and $\mb{u}$ gives
\begin{equation}
\begin{aligned}
m\dot{\mb{v}} &= m\mb{g} + \mb{R}(\sum{H}_i + \begin{bmatrix} 0 & 0 & -T_{\mathrm{tot}}\end{bmatrix}^\top) \\
\mb{J}\dot{\boldsymbol{\omega}} + (\boldsymbol{\omega} \times \mb{J}\boldsymbol{\omega}) &= \begin{bmatrix}\tau_x & \tau_y & \tau_z \end{bmatrix}^\top + \sum{(\mb{r}_i \times \mb{H}_i)}.
\end{aligned}
\end{equation}
These two equations, along with the quaternion derivative $\dot{\mb{q}} = 1/2 \mb{q} \otimes \begin{bmatrix}0 & \boldsymbol{\omega} \end{bmatrix}$ and the relation $\dot{\mb{p}} = \mb{v}$ allows us to form a closed system of equations of the form $\dot{\mb{x}} = f(\mb{x},\mb{u})$.

\end{document}